\section{Triển khai}

Quá trình hiện thực hệ thống RAG được tiến hành trên môi trường Google Colab. Do hạn chế về tài nguyên phần cứng (RAM và VRAM) khi phải đồng thời tải cả Vector Database, Embedding Model và LLM, đề xuất chia quá trình triển khai thành hai giai đoạn độc lập: 
\begin{itemize}
    \item Xây dựng dữ liệu và Truy xuất ngữ cảnh (Context Retrieval).
    \item Sinh câu trả lời (Generation).
\end{itemize}

\subsection{Xây dựng Vector Database và cấu hình Embedding}

Giai đoạn đầu tiên là chuyển đổi dữ liệu văn bản pháp luật thành các vector đặc trưng. Sử dụng
thư viện LangChain và FAISS để quản lý và lập chỉ mục.
 
\begin{itemize}
    \item \textbf{Xử lý dữ liệu:} Các văn bản được chia nhỏ (chunking) với kích thước \texttt{chunk\_size=1024} để
    tối ưu hóa ngữ nghĩa theo kết quả thực nghiệm ở phần trước.
 
    \item \textbf{Embedding Model:} Sử dụng mô hình \texttt{keepitreal/vietnamese-sbert} với cấu hình chuẩn hóa
    vector để phục vụ cho phép đo Cosine Similarity.
 
    \item \textbf{Vector Store:} FAISS được khởi tạo với chiến lược khoảng cách là \texttt{DistanceStrategy.COSINE}
    thay vì $L2$ (Euclidean) mặc định để tăng độ chính xác trong không gian ngữ nghĩa.
\end{itemize}

\subsection{Cấu hình Embedding và FAISS}

Đoạn mã cấu hình Embedding và Vector Store FAISS được trình bày như sau:
 
\begin{verbatim}
# Cấu hình chuẩn hóa vector để sử dụng Cosine Similarity
encode_kwargs = {
    "normalize_embeddings": True
}
 
embeddings = HuggingFaceEmbeddings(
    model_name="keepitreal/vietnamese-sbert",
    model_kwargs={"device": "cuda"},
    encode_kwargs=encode_kwargs
)
 
# Khởi tạo Vectorstore với Cosine Strategy
vectorstore = FAISS.from_documents(
    documents=batch_docs,
    embedding=embeddings,
    distance_strategy=DistanceStrategy.COSINE
)
\end{verbatim}

\subsection{Chiến lược truy xuất Hybrid và tối ưu hóa tài nguyên}
 
Đây là bước quan trọng nhất để giải quyết bài toán giới hạn bộ nhớ. Thay vì truy xuất trực
tiếp trong lúc sinh câu trả lời (Online Retrieval), thực hiện \textit{Pre-computation} (tính toán trước)
ngữ cảnh cho tập dữ liệu kiểm thử (RES Dataset).
 
\subsubsection{Cơ chế Hybrid Search kết hợp Rerank}
 
Xây dựng lớp \texttt{HybridFAISSBM25Searcher} kết hợp giữa tìm kiếm ngữ nghĩa (Dense Retrieval với
FAISS) và tìm kiếm từ khóa (Lexical Retrieval với BM25). Kết quả sau đó được tinh chỉnh lại
thứ hạng (Rerank) bằng mô hình \texttt{BAAI/bge-reranker-large}.
 
\begin{verbatim}
class HybridFAISSBM25Searcher:
 
    def _hybrid_search(self, query, top_k, dense_k, alpha):
        # Dense Retrieval với FAISS
        dense_results = self.faiss_db.similarity_search_with_score(
            query,
            k=dense_k
        )
 
        # Lexical Retrieval với BM25
        bm25_scores = self.bm25.get_scores(
            self.fast_tokenize(query)
        )
 
        # Tính điểm tổng hợp:
        # score = alpha * dense + (1 - alpha) * lexical
        # ... (logic tính toán score)
 
        return results[:top_k]
 
    def _rerank(self, query, docs, top_k):
        # Sử dụng Cross-Encoder để xếp hạng lại
        pairs = [
            (query, doc.page_content)
            for doc in docs
        ]
        scores = self.reranker.predict(pairs)
 
        # ... (logic sắp xếp theo score)
        return ranked[:top_k]
\end{verbatim}
 
\subsubsection{Lưu trữ ngữ cảnh (Context Caching)}
 
Để tránh lỗi \textit{Out of Memory} (OOM) khi tải mô hình ngôn ngữ lớn, thực hiện truy xuất
ngữ cảnh cho toàn bộ câu hỏi trong tập \texttt{res.csv} trước và lưu kết quả vào file trung gian
\texttt{res\_with\_context.csv}.
 
File này chứa các cột ngữ cảnh tương ứng với các chiến lược khác nhau:
 
\begin{itemize}
    \item \texttt{context\_baseline}: Ngữ cảnh từ truy xuất cơ bản.
    \item \texttt{context\_hybrid}: Ngữ cảnh từ FAISS kết hợp BM25.
    \item \texttt{context\_hybrid\_rerank}: Ngữ cảnh sau khi qua bước Reranker.
\end{itemize}
 
Đoạn mã thực hiện lưu trữ ngữ cảnh được trình bày như sau:
 
\begin{verbatim}
contexts_rerank = []
 
# Duyệt qua từng câu hỏi và xây dựng ngữ cảnh
for question in tqdm(
        df["Câu Hỏi"],
        desc="Building contexts"
    ):
 
    hybrid_rerank_results = searcher.search(
        query=question,
        use_rerank=True,
        rerank_k=10
    )
 
    context_rerank = build_context(hybrid_rerank_results)
    contexts_rerank.append(context_rerank)
 
# Lưu file kết quả trung gian
df["context_hybrid_rerank"] = contexts_rerank
df.to_csv("res_with_context.csv", index=False)
\end{verbatim}

\subsection{Giai đoạn Sinh (Generation) với LLM}
 
Sau khi đã giải phóng bộ nhớ khỏi Vector Store và Embedding Model, tiến hành tải mô hình
ngôn ngữ lớn \texttt{Qwen/Qwen3-1.7B} để sinh câu trả lời dựa trên các cột ngữ cảnh đã được chuẩn bị
sẵn trong file CSV.
 
Quy trình thực hiện bao gồm các bước sau:
 
\begin{enumerate}
    \item Tải file \texttt{res\_with\_context.csv}.
    \item Khởi tạo LLM Pipeline với thư viện \texttt{transformers} và \texttt{langchain\_huggingface}.
    \item Xây dựng Prompt Template nhằm ép buộc mô hình chỉ trả lời dựa trên ngữ cảnh được
    cung cấp (Grounding).
\end{enumerate}
 
Đoạn mã sinh câu trả lời được trình bày như sau:
 
\begin{verbatim}
# Khởi tạo LLM
llm = HuggingFacePipeline.from_model_id(
    model_id="Qwen/Qwen3-1.7B",
    task="text-generation",
    device=0, # Sử dụng GPU
    model_kwargs={
        "temperature": 0,
        "max_length": 1024
    }
)
 
def rag_query(question, context):
    prompt = f"""
Bạn là một trợ lý thông minh.
Trả lời câu hỏi dựa trên ngữ cảnh sau.
Nếu không đủ thông tin, hãy nói rõ điều đó,
không được tự bịa đặt.
 
Ngữ cảnh:
{context}
 
Câu hỏi: {question}
 
Trả lời ngắn gọn và chính xác nhất có thể:
"""
    output = llm.generate(
        [prompt],
        max_new_tokens=256
    )
    return output.generations[0][0].text
\end{verbatim}